\documentclass[../main.tex]{subfiles}

\begin{document}

\section{第 8 课：完成一个小型技术报告}

\subsection*{本课目标}

\begin{itemize}
  \item 把前 7 课的内容组合起来。
  \item 能读懂并修改一份简短技术报告模板。
  \item 独立完成标题、章节、公式、图、表、引用和参考文献。
\end{itemize}

\subsection*{最小可运行示例}

\begin{CodeBlock}
\documentclass[12pt,a4paper]{ctexart}
\usepackage[margin=2.2cm]{geometry}
\usepackage{amsmath}
\usepackage{graphicx}
\usepackage{hyperref}
\usepackage[backend=biber,style=numeric]{biblatex}
\addbibresource{bib/references.bib}

\title{用 LaTeX 记录一次简单实验}
\author{你的名字}
\date{\today}

\begin{document}
\maketitle
\section{背景}
这是背景介绍。

\section{方法}
核心公式为
\[
y = kx + b
\]

\section{结论}
参考文献示例：\cite{lamport1994latex}。

\printbibliography
\end{document}
\end{CodeBlock}

\subsection*{简明中文解释}

\begin{enumerate}
  \item 这份例子不是为了教新命令，而是为了把前面学过的命令组合起来。
  \item 你会同时看到标题信息、章节、公式、文献引用等多个部分。
  \item 真正的完整模板已经放在 `\texttt{examples/final\_report\_template.tex}`，里面还额外加入了图和表。
  \item 最终目标不是“抄一遍模板”，而是把模板改成你自己的报告。
\end{enumerate}

\subsection*{改什么、看什么}

打开 `\texttt{examples/final\_report\_template.tex}` 后，建议按下面顺序观察：

\begin{itemize}
  \item 先只改标题、作者、摘要，确认整体框架没问题。
  \item 再逐节替换内容，每替换一节就编译一次。
  \item 最后再处理图、表和参考文献。
\end{itemize}

\subsection*{动手修改任务}

\begin{enumerate}
  \item 把模板标题改成你自己的主题，例如某次实验、某个工具或某次课程作业。
  \item 把“方法”一节里的公式换成和你主题更相关的内容。
  \item 把模板里的图题、表题、引用都改成与你的内容匹配的版本。
  \item 对照 `exercises/final_report_task.md` 完成自检。
\end{enumerate}

\subsection*{常见报错或易错点}

\begin{itemize}
  \item 内容一多后，最常见的问题是少一个花括号、漏一个 `\end{...}` 或标签名没同步修改。
  \item 如果图和表的正文引用显示成 `??`，请再次编译。
  \item 如果参考文献没出现，先检查是否真的写了 `\cite{...}`，再检查编译轮数。
\end{itemize}

\subsection*{对应文件}

本课建议直接修改：`\texttt{examples/final\_report\_template.tex}`，并配合 `\texttt{exercises/final\_report\_task.md}` 完成综合练习。
\end{document}
